\section{Conclusion}\label{sec:conclusion}


This paper studies whether countries with institutionalized trade relationships were more resistant to the impact of 2018 US tariffs on import prices and volumes. Using monthly Census trade data at the product-country level, I construct an event-study framework with a triple-differences interpretation that identifies the differential effect of trade agreements on prices and volumes. I first restrict the sample to tariffs that were applied universally, thus removing the waves of tariffs imposed on China. I compare the restricted sample's aggregate results to those from \textcite{amiti_whos_2020}, whose paper I extend, before implementing the triple differences model and examining heterogeneities across steel and non-steel products.  

The paper finds that trade agreements appear to confer some resilience to unilateral tariff shocks. On aggregate, countries that have a bilateral trade agreement with the US exhibit high initial tariff pass-through to prices followed by partial absorption, while non-agreement countries display sustained complete pass-through. However, pre-trend violations within product-category subsamples preclude causal interpretation of the differential price elasticity. The volume results are stronger: agreement countries gain import share following tariff implementation while non-agreement countries' volumes decline, with the differential satisfying parallel trends and permitting causal interpretation. This aggregate volume effect is driven primarily by steel products, where agreement countries' import volumes rise persistently. Non-steel agreement countries experience only a transient volume gain that converges back to non-agreement levels after roughly six months, suggesting that the resilience-enhancing effect of trade agreements may operate differently across product categories. The evidence is consistent with the hypothesis that stronger institutionalized trade relationships moderate the impact of tariff shocks, though the mechanism appears to be product-specific and the price-side evidence remains descriptive rather than causal.